% GENERATED FILE. Edit canonical lessons, then run npm run generate:books.
% canonical-source-hash: fnv1a64:829a98cb313c0789
% canonical-lessons: BN-C400-thikana, BN-C400-chihno, BN-C400-ulto, BN-C400-nirdesh, BN-C400-inggit-kora

\chapter{Addresses and Signs}
\label{ch:bn-addresses-and-signs}
\hlchaptermodality{400}

\begin{chapteropening}
\textbf{By the end of this chapter:} \emph{I can follow an address, a sign and an instruction, and point the way.}

Complete the last lesson of chapter 400: I can follow an address, a sign and an instruction, and point the way.
\end{chapteropening}

\section[\texorpdfstring{\bn{ঠিকানা} (ṭhikānā)}{ঠিকানা (ṭhikānā)}]{\bn{ঠিকানা} (ṭhikānā) — an address}

\label{lesson:BN-C400-thikana}



\begin{quote}
Before the new one: say the Bengali for to warm oneself, then the Bengali for a guide (for visitors).
\end{quote}

\subsection*{You'll want to know: \bn{ঠিকানা}}
\textbf{\bn{ঠিকানা}} — \emph{ṭhikānā} — \textquotedblleft{}an address\textquotedblright{}.

\begin{quote}
\bn{আমি} \bn{বন্ধুর} \bn{নতুন} \bn{ঠিকানা} \bn{খুঁজে} \bn{পেলাম।}
\end{quote}

\begin{quote}
\emph{āmi bôndhur nôtun ṭhikānā khũje pelām.}
\end{quote}

\begin{quote}
I found my friend's new address.
\end{quote}

It shares its first syllable with \bn{ঠিক}, right: an address is the place a thing is fixed.

\begin{cousinweb}
\textbf{\bn{ঠিকানা}} (\emph{ṭhikānā}) is inherited: CDIAL lists Bengali \emph{ṭhikānā}, \textquotedblleft{}a station, a fixed place\textquotedblright{}, beside \emph{ṭhik}, \textquotedblleft{}correct\textquotedblright{}, under the reconstructed \emph{ṭhīkka}, \textquotedblleft{}firm, right\textquotedblright{} (CDIAL 5503); the sense \textquotedblleft{}a postal address\textquotedblright{} is the modern use.
\end{cousinweb}

\begin{grammarlens}[title={What you've built}]
One more word for getting around.
\end{grammarlens}

\subsection*{Your turn}
\begin{itemize}
\raggedright
  \item \emph{Say it:} \emph{ṭhikānā}
  \item \emph{Say it:} \emph{ṭhikānā}, once more
  \item \emph{Say it:} \emph{āmi bôndhur nôtun ṭhikānā khũje pelām.} — I found my friend's new address.
  \item \emph{Recall:} say the Bengali for to warm oneself, then the Bengali for a guide (for visitors), then say \emph{ṭhikānā} again
\end{itemize}

\begin{tcolorbox}[breakable,colback=teal!4,colframe=teal!35!black,title={Before you move on}]
What does \bn{ঠিকানা} mean? (\textquotedblleft{}An address\textquotedblright{}.) Say it once more.
\end{tcolorbox}
\section[\texorpdfstring{\bn{চিহ্ন} (chihnô)}{চিহ্ন (chihnô)}]{\bn{চিহ্ন} (chihnô) — a sign; a mark}

\label{lesson:BN-C400-chihno}



\begin{quote}
Before the new one: say the Bengali for a guide (for visitors), then the Bengali for an address.
\end{quote}

\subsection*{You'll want to know: \bn{চিহ্ন}}
\textbf{\bn{চিহ্ন}} — \emph{chihnô} — \textquotedblleft{}a sign; a mark\textquotedblright{}.

\begin{quote}
\bn{রাস্তার} \bn{চিহ্ন} \bn{দেখে} \bn{চালক} \bn{ধীরে} \bn{চালালেন।}
\end{quote}

\begin{quote}
\emph{rāstār chihnô dekhe chālôk dhīre chālālen.}
\end{quote}

\begin{quote}
Seeing the road sign, the driver drove slowly.
\end{quote}

Said chinho: the h is heard before the n.

\begin{cousinweb}
\textbf{\bn{চিহ্ন}} (\emph{chihnô}) is Sanskrit \emph{cihna}, \textquotedblleft{}a mark, a stamp, a sign\textquotedblright{} (Monier-Williams), taken over whole.
\end{cousinweb}

\begin{grammarlens}[title={What you've built}]
One more word for getting around.
\end{grammarlens}

\subsection*{Your turn}
\begin{itemize}
\raggedright
  \item \emph{Say it:} \emph{chihnô}
  \item \emph{Say it:} \emph{chihnô}, once more
  \item \emph{Say it:} \emph{rāstār chihnô dekhe chālôk dhīre chālālen.} — Seeing the road sign, the driver drove slowly.
  \item \emph{Recall:} say the Bengali for a guide (for visitors), then the Bengali for an address, then say \emph{chihnô} again
\end{itemize}

\begin{tcolorbox}[breakable,colback=teal!4,colframe=teal!35!black,title={Before you move on}]
What does \bn{চিহ্ন} mean? (\textquotedblleft{}A sign; a mark\textquotedblright{}.) Say it once more.
\end{tcolorbox}
\section[\texorpdfstring{\bn{উল্টো} (ulṭo)}{উল্টো (ulṭo)}]{\bn{উল্টো} (ulṭo) — opposite; the other way; reversed}

\label{lesson:BN-C400-ulto}



\begin{quote}
Before the new one: say the Bengali for an address, then the Bengali for a sign.
\end{quote}

\subsection*{You'll want to know: \bn{উল্টো}}
\textbf{\bn{উল্টো}} — \emph{ulṭo} — \textquotedblleft{}opposite; the other way; reversed\textquotedblright{}.

\begin{quote}
\bn{আমরা} \bn{ভুল} \bn{করে} \bn{উল্টো} \bn{দিকের} \bn{বাসে} \bn{উঠলাম।}
\end{quote}

\begin{quote}
\emph{āmrā bhul kore ulṭo diker bāse uṭhlām.}
\end{quote}

\begin{quote}
By mistake we got on the bus going the other way.
\end{quote}

\bn{উল্টো} \bn{দিকে}, in the opposite direction; \bn{উল্টো} \bn{দিকের} \bn{দোকান}, the shop across the road.

\begin{cousinweb}
\textbf{\bn{উল্টো}} (\emph{ulṭo}) is inherited: CDIAL lists Bengali \emph{ulṭo}, \textquotedblleft{}reversed\textquotedblright{}, under the reconstructed \emph{ullaṭati}, \textquotedblleft{}turns over\textquotedblright{} (CDIAL 2368), the same root as \bn{উল্টানো}, to overturn.
\end{cousinweb}

\begin{grammarlens}[title={What you've built}]
One more word for getting around.
\end{grammarlens}

\subsection*{Your turn}
\begin{itemize}
\raggedright
  \item \emph{Say it:} \emph{ulṭo}
  \item \emph{Say it:} \emph{ulṭo}, once more
  \item \emph{Say it:} \emph{āmrā bhul kore ulṭo diker bāse uṭhlām.} — By mistake we got on the bus going the other way.
  \item \emph{Recall:} say the Bengali for an address, then the Bengali for a sign, then say \emph{ulṭo} again
\end{itemize}

\begin{tcolorbox}[breakable,colback=teal!4,colframe=teal!35!black,title={Before you move on}]
What does \bn{উল্টো} mean? (\textquotedblleft{}Opposite; the other way; reversed\textquotedblright{}.) Say it once more.
\end{tcolorbox}
\section[\texorpdfstring{\bn{নির্দেশ} (nirdesh)}{নির্দেশ (nirdesh)}]{\bn{নির্দেশ} (nirdesh) — an instruction; directions}

\label{lesson:BN-C400-nirdesh}



\begin{quote}
Before the new one: say the Bengali for a sign, then the Bengali for opposite.
\end{quote}

\subsection*{You'll want to know: \bn{নির্দেশ}}
\textbf{\bn{নির্দেশ}} — \emph{nirdesh} — \textquotedblleft{}an instruction; directions\textquotedblright{}.

\begin{quote}
\bn{পুলিশের} \bn{নির্দেশে} \bn{আমরা} \bn{গাড়ি} \bn{থামালাম।}
\end{quote}

\begin{quote}
\emph{pulisher nirdeshe āmrā gāṛi thāmālām.}
\end{quote}

\begin{quote}
On the police's instruction we stopped the car.
\end{quote}

\begin{cousinweb}
\textbf{\bn{নির্দেশ}} (\emph{nirdesh}) is Sanskrit \emph{nirdeśa}, \textquotedblleft{}pointing out, directing, an order\textquotedblright{} (Monier-Williams), from \emph{diś}, to point.
\end{cousinweb}

\begin{grammarlens}[title={What you've built}]
One more word for getting around.
\end{grammarlens}

\subsection*{Your turn}
\begin{itemize}
\raggedright
  \item \emph{Say it:} \emph{nirdesh}
  \item \emph{Say it:} \emph{nirdesh}, once more
  \item \emph{Say it:} \emph{pulisher nirdeshe āmrā gāṛi thāmālām.} — On the police's instruction we stopped the car.
  \item \emph{Recall:} say the Bengali for a sign, then the Bengali for opposite, then say \emph{nirdesh} again
\end{itemize}

\begin{tcolorbox}[breakable,colback=teal!4,colframe=teal!35!black,title={Before you move on}]
What does \bn{নির্দেশ} mean? (\textquotedblleft{}An instruction; directions\textquotedblright{}.) Say it once more.
\end{tcolorbox}
\section[\texorpdfstring{\bn{ইঙ্গিত} \bn{করা} (inggit kôrā)}{ইঙ্গিত করা (inggit kôrā)}]{\bn{ইঙ্গিত} \bn{করা} (inggit kôrā) — to point (to); to hint at}

\label{lesson:BN-C400-inggit-kora}



\begin{quote}
Before the new one: say the Bengali for opposite, then the Bengali for an instruction.
\end{quote}

\subsection*{You'll want to know: \bn{ইঙ্গিত} \bn{করা}}
\textbf{\bn{ইঙ্গিত} \bn{করা}} — \emph{inggit kôrā} — \textquotedblleft{}to point (to); to hint at\textquotedblright{}.

\begin{quote}
\bn{দোকানদার} \bn{ডানদিকের} \bn{গলির} \bn{দিকে} \bn{ইঙ্গিত} \bn{করলেন।}
\end{quote}

\begin{quote}
\emph{dokāndār ḍāndiker gôlir dike inggit kôrlen.}
\end{quote}

\begin{quote}
The shopkeeper pointed towards the lane on the right.
\end{quote}

Past (he/she …): \bn{তিনি} \bn{নদীর} \bn{দিকে} \bn{ইঙ্গিত} \bn{করলেন}, \emph{tini nôdīr dike inggit kôrlen}.

\bn{ইঙ্গিত} alone is a sign or a hint: \bn{চোখের} \bn{ইঙ্গিত}, a sign with the eyes.

\begin{cousinweb}
\textbf{\bn{ইঙ্গিত} \bn{করা}} (\emph{inggit kôrā}) joins Sanskrit \emph{iṅgita}, \textquotedblleft{}a motion of the body that shows what one means, a hint\textquotedblright{} (Monier-Williams), to \bn{করা}, to do; \bn{ঙ্গ} is said ngg.
\end{cousinweb}

\begin{grammarlens}[title={What you've built}]
One more word for getting around.
\end{grammarlens}

\subsection*{Your turn}
\begin{itemize}
\raggedright
  \item \emph{Say it:} \emph{inggit kôrā}
  \item \emph{Say it:} \emph{inggit kôrā}, once more
  \item \emph{Say it:} \emph{dokāndār ḍāndiker gôlir dike inggit kôrlen.} — The shopkeeper pointed towards the lane on the right.
  \item \emph{Recall:} say the Bengali for opposite, then the Bengali for an instruction, then say \emph{inggit kôrā} again
\end{itemize}

\begin{tcolorbox}[breakable,colback=teal!4,colframe=teal!35!black,title={Before you move on}]
What does \bn{ইঙ্গিত} \bn{করা} mean? (\textquotedblleft{}To point (to); to hint at\textquotedblright{}.) Say it once more.
\end{tcolorbox}
